\documentclass[10pt,a4paper]{amsart}
\usepackage[margin=19mm]{geometry}
\usepackage{amsmath,amssymb,mathtools}
\usepackage[scr=rsfs,cal=euler]{mathalfa}
\usepackage{xfrac}
\usepackage[unicode,hidelinks,bookmarks=false]{hyperref}
\newcommand{\ie}{i.e.\ }
\newcommand{\cf}{cf.\ }
\newcommand{\resp}{respectively\ }
\newcommand{\Subsection}{\subsection}
\providecommand{\nobreakdash}{\nobreak\hbox{-}}
\setlength{\emergencystretch}{2em}
\multlinegap0pt
\usepackage{amssymb}
\usepackage{mathtools}
\usepackage[shortlabels]{enumitem}
\newtheorem{thm}{Theorem}
\newtheorem{lemma}{Lemma}
\usepackage{cleveref}
\crefname{thm}{Theorem}{Theorems}
\crefname{lemma}{Lemma}{Lemmas}
\DeclareMathOperator{\BL}{BL}
\DeclareMathOperator{\Ell}{\mathscr{L}}
\DeclareMathOperator{\End}{End}
\DeclareMathOperator{\Gr}{Gr}
\newcommand{\R}{\mathbb{R}}
\DeclareMathOperator{\Span}{Span}
\newcommand{\bR}{\mathbf{R}}
\newcommand{\eps}{\varepsilon}
\DeclareMathOperator{\id}{id}
\newcommand{\norm}[1]{\lVert#1\rVert}   % norme
\DeclareMathOperator{\rk}{rk}
\newcommand{\sD} {{\mathscr D}}
\newcommand{\sN}{\mathscr{N}}
\newcommand{\setBig}[1]{\Bigl\{\, #1 \,\Bigr\}}
\renewcommand{\subset}{\subseteq}
\title{Effective Brascamp-Lieb inequalities}
\author{Timoth\'ee B\'enard, Weikun He}
\date{Source: arXiv:2511.11091v2 (CC BY 4.0)}
\begin{document}
\maketitle

\noindent\emph{Source and licence.} Effective Brascamp-Lieb inequalities. Version arXiv:2511.11091v2 (CC BY 4.0), distributed by arXiv under the Creative Commons Attribution 4.0 International licence, http://creativecommons.org/licenses/by/4.0/, as recorded in the arXiv licence metadata for this exact version and shown on the abstract page https://arxiv.org/abs/2511.11091v2. Full licence notice: \texttt{licenses/arxiv-2511.11091-CC-BY-4.0.txt}.\par\smallskip

\noindent\textit{Editorial scope.} Counted unit: the article's thm thm:loc-lb (source0.tex lines 729--737). The article's complete original proof of it is delivered with it (source0.tex lines 1240--1278, one \texttt{proof} environment of 39 lines, which the article places away from the statement). No mathematics is written, repaired or re-derived: the statement and the proof are the article's own lines, in the article's own notation, and no retained line is substituted or re-flowed.  Retained uncounted as the source-local context the proof needs: lines 643--645; lines 869--876; lines 1213--1219.  Nothing was omitted from any retained span, so the package carries no omission record.  No retained line contains a citation, so the wrapper carries no bibliography and no bibliography entry.  No retained line contains a figure environment, an image-inclusion command or drawing code, so no image is delivered and the package declares no asset dependency.  Editorial formatting only, in addition: an AMS \texttt{amsart} preamble that restates the article's own declarations and macros used by the retained lines, namely the environments \texttt{thm}, \texttt{lemma} on separate counters, together with the standard packages those lines need.  The retained environments print in this document as Theorem 1, Lemma 1 in order of appearance, whereas the article prints the same items with the numbering of its own full document.  The complete native source of the article as distributed in the arXiv e-print for this exact version is bundled unchanged as \texttt{sources/189092/source0.tex}.
\begin{equation}\label{def-chi-sN}
\chi_{R}(x)= e^{- \pi \langle x, R x \rangle } \quad \text{ and }\quad \sN_{R}(x)=(\det R)^{1/2}\chi_{R}(x).
\end{equation}

\begin{lemma}
\label{lm:arank}
For $\alpha \in \R_{\geq 0}$, $\ell \in \Ell(H,H')$,  $W \in \Gr(H)$, we have
\[
\rk_{\alpha}(\ell \mid W) = \min \setBig{\dim E \,:\, E \in \Gr(H') \text{ with } \ell \bigl(B^W_1\bigr) \subseteq B^{H'}_\alpha + E}.
\]
%Moreover, considering a Cartan basis $(w_{i})_{i=1}^k$ of $W$ adapted to $\ell_{|W}$, we may choose $E$ to be  $\Span \{w_{i}\,:\, \|\ell(w_{i})\|>\alpha\}$.
\end{lemma}

\begin{lemma}\label{eq:Choleski}
Let $H$ be a Euclidean space, let $Q\in \End(H)$ be a positive semi-definite symmetric endomorphism,  let  $(\eps_{1}, \dots, \eps_{d})$ be a basis of $H$.
Then
\begin{equation*}
\det Q \leq \|\eps_{1}\wedge \dots \wedge \eps_{d}\|^{-2}\prod_{k=1}^d \langle Q \eps_{k}, \eps_{k}\rangle.
\end{equation*}
\end{lemma}

\begin{thm}[Lower bound]\label{thm:loc-lb}
Let $\sD = \bigl( (\ell_j)_j, (q_j)_j \bigr)$ be a Brascamp-Lieb datum. Set $C:=1+\sup_{j\in J} \|\ell_{j}\|$.
Let $\bR = (\id_{H_j})_{j \in J}$ and $T = t \id_H$ for some $t>0$.

Then for all $\alpha \in (0, 1]$ and all $W \in \Gr(H)$, we have
\[
\BL( \sD, \bR, T) \geq \left( (C/\alpha)^2 t+ C^2\sum\nolimits_{j \in J} q_j \right)^{-\dim H / 2} \alpha^{\sum_{j \in J} q_j \rk_{\alpha}(\ell_j \mid W) - \dim W}.
\]
\end{thm}

\begin{proof}[Proof of \Cref{thm:loc-lb}]
Consider a collection  $(A_j)_{j \in J}$ of positive definite symmetric endomorphisms $A_{j}\in \End(H_{j})$ with $A_{j}\leq  \id_{H_j}$. The localised regularised  Brascamp-Lieb inequality with  test functions $f_{j}=\chi_{A_{j}}$ (definition in \eqref{def-chi-sN}) yields by direct computation
\begin{equation}\label{BLDRT-lwb}
\BL(\sD, \bR, T) \geq \left(\frac{\prod_{j \in J} (\det A_j)^{q_{j}}}{\det\left(T + \sum\nolimits_{j \in J} q_j \ell_j^* A_j \ell_j \right)}\right)^{1/2}.
\end{equation}
The strategy is to choose the $A_{j}$ so that the right-hand side of \eqref{BLDRT-lwb} has the desired lower bound.


By \Cref{lm:arank}, for each $j \in J$, there exists $E_j \in \Gr(H_j)$ of dimension $\dim E_j = \rk_\alpha(\ell_j \mid W)$ such that
\[
\ell_j B^W_1 \subset B^{H_j}_\alpha + E_j.
\]
Let $A_j \in \End(H_j)$ be the positive definite symmetric endomorphism whose square root is
\[
A_j^{1/2} = \alpha \id_{E_j} \oplus \id_{E_j^\perp}.
\]
Observe that $A_j \leq \id_{H_j}$ and $\det(A_j)^{1/2} = \alpha^{\rk_\alpha(\ell_j \mid W)}$. It remains to bound from above the determinant of $T+\sum_{j \in J} q_j \ell_j^* A_j \ell_j$. We assume from the start that $E_{j}$ has the property described in the remark after \Cref{lm:arank}.
In particular, for any unit vector $w \in W$, we can decompose $w = w' + w''$ into orthogonal vectors satisfying $\ell_j w' \in E_j$, $\ell_j w'' \in E_j^\perp$, $\norm{\ell_j w'} \leq C\norm{w'}$ and $\norm{\ell_j w''} \leq \alpha \norm{w''}$.
It follows that
\begin{align*}
\langle \ell_j^* A_j \ell_j w, w \rangle &= \norm{A_j^{1/2} \ell_j w}^2 \\
&= \norm{ \alpha\ell_j w' + \ell_j w''}^2\\
&= \alpha^2 \norm{\ell_j w'}^2 + \norm{\ell_j w''}^2\\
&\leq \alpha^2 C^2\norm{w'}^2 + \alpha^2 \norm{w''}^2\\
&\leq \alpha^2 C^2
\end{align*}
On the other hand, for every unit vector $w_{\perp} \in W^\perp$, we have
\[
\langle \ell_j^* A_j \ell_j w_{\perp}, w_{\perp} \rangle = \norm{A_j^{1/2} \ell_j w_{\perp}}^2 \leq C^2.
\]
Consider now an orthonormal basis $(e_1, \dotsc, e_d)$ of $H$, obtain by joining an orthonormal basis of $W$ with one of $W^\perp$.
Using \Cref{eq:Choleski}, the above inequalities, then $C\geq 1\geq \alpha$, we find
\begin{align*}
\det\left(T + \sum\nolimits_{j \in J} q_j \ell_j^* A_j \ell_j \right) &\leq \prod_{i = 1}^d \left\langle Te_i + \sum\nolimits_{j \in J} q_j \ell_j^* A_j \ell_j e_i, \,e_i \right\rangle\\
&\leq \Big(t + \sum\nolimits_{j \in J} q_j \alpha^2C^2 \Bigr)^{\dim W} \left(t + \sum\nolimits_{j \in J} q_j C^2\right)^{\dim W^\perp}\\
&\leq  \alpha^{2 \dim W}\left( (C/\alpha)^2 t + C^2\sum\nolimits_{j \in J} q_j \right)^d.
\end{align*}
By combination with \eqref{BLDRT-lwb}, this yields the announced lower bound.
\end{proof}

\end{document}
